\documentclass[letterpaper]{article} % DO NOT CHANGE THIS
\usepackage[submission]{aaai2027}  % DO NOT CHANGE THIS
\usepackage[hyphens]{url}  % DO NOT CHANGE THIS
\usepackage{graphicx} % DO NOT CHANGE THIS
\urlstyle{rm} % DO NOT CHANGE THIS
\def\UrlFont{\rm}  % DO NOT CHANGE THIS
\usepackage{natbib}  % DO NOT CHANGE THIS AND DO NOT ADD ANY OPTIONS TO IT
\usepackage{caption} % DO NOT CHANGE THIS AND DO NOT ADD ANY OPTIONS TO IT
\frenchspacing  % DO NOT CHANGE THIS
\usepackage{booktabs}
\usepackage{amsmath}
\pdfinfo{
/TemplateVersion (2027.1)
}
\setcounter{secnumdepth}{0}

\title{The Missing Band Is a Distribution:\\Do Generative Bandwidth-Extension Models Sample It?}
\author{Anonymous submission}
\affiliations{}

\begin{document}
\maketitle

\begin{abstract}
Generative audio bandwidth extension (BWE) models are presented as samplers of the missing high band, but
are evaluated with single-reference metrics that cannot tell a sampler from a regressor. We audit four
released 12$\to$48\,kHz systems with the checks used for stochastic super-resolution of weather fields (fair
CRPS, rank histograms, spread), after showing that when the high band is incoherent with the truth,
waveform spread measures output level as much as diversity. FLowHigh ranks first on LSD and ViSQOL yet is a
regressor: one Euler step returns the conditional mean for any start noise. NU-Wave~2 samples, but
8--9\,dB too quietly; AudioSR is dispersed but biased by band; every model pulls each utterance's level
toward the corpus mean. On one 88k-parameter network, the deficit is set mostly by the loss domain and is
closed by an energy score at no inference cost.
\end{abstract}

\section{Introduction}
Many high bands are consistent with one narrow band, which bounds what any deterministic BWE estimator can
recover \citep{jax2002,bruna2016}. A point estimate trained on the waveform returns a conditional mean, close
to silence for an incoherent band: a muffled output.
Diffusion and flow models \citep{han2022,liu2024,yun2025} are meant to avoid this by sampling $p(y\mid x)$,
yet they are scored by comparing one output with one reference, which rewards an empty band \citep{su2021}.
We found no BWE work that checks whether these models sample. Stochastic downscaling of weather fields,
super-resolution by another name, checks this with proper scores \citep{gneiting2007} and rank histograms
\citep{leinonen2021}. We bring those checks to BWE, with (i)~a separation of level from dispersion that keeps
them meaningful when the band is incoherent, (ii)~an audit of four released systems through one scorer, and
(iii)~a controlled study separating the loss domain from sampling.

\section{Setup}
VCTK~0.92, test split of FLowHigh and NU-Wave~2: eight speakers unseen in training by every model except
possibly AudioSR, with their input protocol (order-8 Chebyshev low-pass, 12\,kHz). Levels, LSD and ViSQOL (v3,
audio mode) use 240 utterances; calibration uses 32 utterances $\times$ 8 draws. Outputs are scaled to the
truth's energy below 5.5\,kHz. The \emph{deficit} is $10\log_{10}$ of output over true energy per third octave
above 6\,kHz, averaged over bands. CRPS is the fair ensemble estimator on high-band (HB) log-magnitude; a
one-output model scores its MAE. The harness reproduces published LSD/ViSQOL: 0.761/3.61 for FLowHigh
(reported 0.75/3.61), 0.780/3.48 for AP-BWE \citep{lu2024} (0.78/3.46); NU-Wave~2 is slightly worse,
0.989/2.57 against 0.94/2.75.

\section{Level Before Dispersion}
Let the truth's HB have energy $E$, the output's $pE$, and let $c=\mathrm{Re}\langle y_H,\hat y_H\rangle/E$ be
the coherent fraction. With a correct low band, the squared error is
\begin{equation}
\|y-\hat y\|^2 = E\,(1+p-2c).
\end{equation}
An empty band costs $E$, and only $c>p/2$ beats it, so when $c\approx0$ silence maximises SNR. On our
utterances this ceiling is 22.02\,dB at 12\,kHz, and a reported SNR $S$ bounds the level:
$p\le10^{(C-S)/10}-1+2c$. NU-Wave~2's own table ($C-S=0.5$\,dB) gives at most $-8.5$\,dB; we measure a
broadband deficit of $-8.9$\,dB. Every model has $c<0.02$, including FLowHigh with attention over the
whole utterance, so none recovers the HB waveform.

The same identity governs spread. For $M$ draws independent of one another and of the truth, the gap
$10\log_{10}\bigl(\overline{\|y-x_i\|^2}/\|y-\bar x\|^2\bigr)$ is, in expectation,
$10\log_{10}\frac{1+p}{1+p/M}$: a function of level alone. A calibrated sampler has $p=1$ and a gap of 2.50\,dB at $M=8$. We therefore report
level first, and judge diversity by comparing each model's gap with the independent value \emph{at its own
level}.

\begin{table}[t]
\centering\small\setlength{\tabcolsep}{2.4pt}
\begin{tabular}{lrrrrr}
\toprule
model & deficit & LSD & ViSQOL & gap (indep.) & CRPS \\
\midrule
FLowHigh 49M & $-1.60$ & 0.761 & 3.61 & 0.08 (1.61)$^\dagger$ & 0.682$^\dagger$ \\
AP-BWE 30M & $-1.38$ & 0.780 & 3.48 & -- & 0.777$^*$ \\
NU-Wave 2 1.7M & $-7.96$ & 0.989 & 2.57 & 0.57 (0.47) & 0.583 \\
AudioSR 672M & $-0.60$ & 1.561 & 2.41 & 2.62 (3.45) & 1.243 \\
ours, point 88k & $-6.18$ & 0.864 & 3.13 & -- & 0.889$^*$ \\
ours, sampler 88k & $-0.70$ & 0.895 & 2.83 & 2.07 (2.09) & 0.514 \\
\bottomrule
\end{tabular}
\caption{One scorer. Deficit (dB), LSD, ViSQOL: 240 utterances. Gap (dB; independent-draw value at the
model's level) and fair CRPS: 32 utterances $\times$ 8 draws. $^\dagger$Prior restored. $^*$CRPS = MAE.}
\label{tab:released}
\end{table}

\section{What the Audit Finds}
\textbf{FLowHigh is a regressor.} It is trained from $x_0+\epsilon$ (the input's mel spectrogram plus noise)
toward $(x_1-x_0)-(1-\sigma)\epsilon$, and it also receives $x_0$ as conditioning, so at $t=0$ it can recover
$\epsilon$. The optimal field then returns $\mathrm{E}[x_1\mid x_0]+\sigma s\epsilon$ after one Euler step, for any
start-noise scale $s$; the released script also resets $s$ to $10^{-4}$.
With the trained prior restored, its draws remain almost identical: the gap is 0.08\,dB against 1.61 for
independent draws at its level, and the truth lies below or above all eight draws in 41\% and 46\% of HB
bins. LSD and ViSQOL rank it first, and nothing in them registers that it does not sample.

\textbf{NU-Wave~2 samples, quietly.} Its gap (0.57\,dB) matches independent draws at its level (0.47): the
draws are as diverse as its level allows. What fails is the level: 9.0\,dB under on loud frames and 4.7\,dB
over on quiet ones. Its small spread (0.094) is a level problem, not a diversity problem.

\textbf{AudioSR is dispersed but biased by band} ($+5$\,dB at 11\,kHz, $-11$\,dB at 21\,kHz; truth above
every draw in 52\% of HB bins). AP-BWE is a point estimate by design.

\textbf{Level regresses per utterance.} Regressing each utterance's deficit on the truth's HB energy share gives slopes of $-0.07$ (FLowHigh), $-0.12$ (AP-BWE), $-0.32$ (NU-Wave~2) and
$-0.34$ (AudioSR); our arms lie between $-0.20$ and $-0.35$. Bright utterances come out too dark and dark
ones too bright.

\begin{table}[t]
\centering\small\setlength{\tabcolsep}{4pt}
\begin{tabular}{llrrr}
\toprule
arm & objective & deficit & gap & CRPS \\
\midrule
det\_paper & waveform L1 (LISA) & $-21.78$ & -- & 2.853$^*$ \\
det & + log-mag., $\lambda{=}0.01$ & $-6.18$ & -- & 0.889$^*$ \\
es\_marg & energy score & $-1.05$ & 2.21 & 0.537 \\
es\_erb & + ERB bands & $-1.08$ & 2.03 & 0.512 \\
es\_erb & $\lambda{=}0.001$ & $-5.73$ & 1.16 & 0.683 \\
es\_dec\_erb & + decoder noise, $\lambda{=}0.1$ & $-0.70$ & 2.07 & 0.514 \\
\bottomrule
\end{tabular}
\caption{One 88k network, identical batches; columns as in Table~\ref{tab:released}.}
\label{tab:arms}
\end{table}

\section{What Sets the Deficit}
All arms retrain LISA \citep{kim2022} on identical batches for 104{,}950 steps; the samplers add 8 Gaussian
input channels, and 4 per output sample in the decoder-noise arm (Table~\ref{tab:arms}). Its released L1 loss leaves the band
21.8\,dB short. Adding a log-magnitude term to the same point estimate recovers 15.6\,dB; replacing the point
loss by a two-draw energy score \citep{gritsenko2020,bouchacourt2016} recovers 5.1\,dB more. The loss domain, not point estimation as such, does most of the work: FLowHigh, a conditional
mean in the log-mel domain, is not quiet. The ERB term lowers the band-correlation error from 0.374 to 0.289.
Our sampler is not calibrated either: the truth lies above all draws in 21\% of HB bins against an ideal 11\%,
while the lower tail is at its ideal, a residual level bias rather than a narrow ensemble.

\section{Conclusion}
Single-reference metrics cannot tell whether a BWE model samples, and when the band is incoherent, spread
largely measures level. We suggest reporting level (per band and frame loudness), diversity relative to the
model's own level, and rank tails beside LSD and ViSQOL; CRPS favours samplers over point models by
construction, so we compare it among samplers only. Limitations: no listening test, 32 calibration
utterances, one seed, VCTK only.

\bibliography{refs}

\end{document}
